\section{Apphost: Architektur, Umsetzung und Bedrohungsmodell} \label{Apphost}

Dieses Kapitel beschreibt den zentralen Baustein des Gesamtsystems, den
sogenannten \emph{Apphost}. Auf ihm laufen alle Alltags- und Produktivdienste
des Mandanten. Zunächst werden die Entscheidungen erläutert, die zu der
gewählten Architektur geführt haben. Anschließend wird der konkrete Aufbau
beschrieben. Abschließend wird gezeigt, wie die Umsetzung das vorgegebene
Bedrohungsmodell abdeckt. An geeigneten Stellen wird auf die separate
Installationsanleitung sowie auf den Anhang verwiesen, in dem die einzelnen
Schritte und alle Sicherheitseinstellungen im Detail aufgeführt sind.

\subsection{Evaluation verschiedener Ansätze}

Vor der Festlegung der endgültigen Architektur wurden mehrere Varianten
gegeneinander abgewogen. Gesucht wurde eine Lösung, die sicher und
reproduzierbar ist und zugleich wartbar bleibt. Die Nachvollziehbarkeit war
dabei das ausschlaggebende Kriterium, denn der Mandant soll die Infrastruktur
selbst prüfen und betreiben können.

\subsubsection*{Verworfen: Kubernetes mit Talos Linux}

Der erste Entwurf basierte auf Kubernetes als Orchestrierung und Talos Linux
als minimalem, \ac{API}-gesteuertem Betriebssystem für die Nodes. Talos
besitzt weder eine Shell noch einen \ac{SSH}-Zugang. Die gesamte Konfiguration
erfolgt über eine deklarative \ac{API}. Kubernetes bringt zudem native
Konzepte wie Network Policies und Secrets mit.

Dieser Ansatz wurde bereits weitgehend umgesetzt. Im Verlauf der Arbeit zeigte
sich jedoch, dass die Komplexität von Kubernetes in keinem angemessenen
Verhältnis zum Ziel der Aufgabe steht. Für einen einzelnen Mandanten mit einer
überschaubaren Zahl privat betriebener Dienste bedeutet ein Kubernetes-Cluster
einen erheblichen dauerhaften Wartungsaufwand. Ingress-Controller,
\ac{CNI}-Plugin, Storage-Provisioner, Zertifikatsverwaltung und
Cluster-Updates müssen verstanden und gepflegt werden. Der Kern der Aufgabe
liegt jedoch in der \emph{sorgfältigen Absicherung} der Dienste und nicht im
Betrieb eines Cluster-Stacks. Aus diesem Grund wurde der Stack trotz des
bereits investierten Aufwands auf Docker Compose neu aufgebaut (siehe
Abschnitt~\ref{sec:herausforderungen-selfhosted}).

\subsubsection*{Ausgetauschte Komponenten}

Mit dem Wechsel weg von Kubernetes wurden auch einige zuvor gewählte Werkzeuge
ersetzt, da sie ohne Cluster keinen Mehrwert boten oder für den Anwendungsfall
überdimensioniert waren:

\begin{itemize}
  \item \textbf{Envoy Gateway zu Traefik.} Envoy Gateway ist eng an die
    Kubernetes Gateway \ac{API} gekoppelt, sodass diese Integration ohne
    Kubernetes entfällt. Traefik liest Container-Labels direkt aus Docker und
    bringt die \ac{ACME}-Zertifikatsverwaltung (Let's Encrypt) einschließlich
    der \ac{DNS}-01-Challenge mit. Dadurch lassen sich gültige Zertifikate
    beziehen, ohne dass der Server aus dem Internet erreichbar sein muss.
  \item \textbf{Keycloak zu Authelia.} Keycloak ist ein vollwertiger
    Identity-Provider für große Organisationen und entsprechend
    schwergewichtig. Authelia lässt sich vollständig deklarativ über
    Konfigurationsdateien einrichten und fügt sich damit besser in den
    Infrastructure-as-Code-Ansatz ein. Für einen einzelnen Nutzer mit wenigen
    Diensten ist es zudem angemessener dimensioniert. Authelia deckt sowohl
    klassisches Forward-Auth als auch \ac{OIDC} ab und erfüllt damit alle
    benötigten Fälle.\cite{keyCloak}
  \item \textbf{Dragonfly zu Valkey.} Ausschlaggebend für Dragonfly war
    ursprünglich die Verfügbarkeit eines Kubernetes-Operators. Valkey steht
    unter dem Dach der Linux Foundation und ist offener lizenziert. Es passt
    damit besser zu einem auf Datenschutz und Unabhängigkeit ausgerichteten
    System. In reinen Benchmarks ist Valkey etwas langsamer als Dragonfly.
    Für den vorliegenden Anwendungsfall, das Caching für Immich und
    Paperless, ist dieser Unterschied ohne Bedeutung.
\end{itemize}

\subsubsection*{Ein einziger gehärteter Apphost statt \ac{VM} pro Dienst}

Eine weitere zentrale Entscheidung war, alle Dienste auf einer einzigen, gut
abgesicherten virtuellen Maschine zu betreiben, statt jeden Dienst in eine
eigene \ac{VM} zu legen. Dafür gab es mehrere Gründe:

\begin{itemize}
  \item \textbf{Geringere Komplexität.} Bei über 18 Diensten wären ebenso
    viele Betriebssysteme einzeln zu aktualisieren, zu härten und zu
    überwachen. Für den Endanwender ist das kaum handhabbar.
  \item \textbf{Kein internes Routing nötig.} Bei vielen \acp{VM} wäre ein
    internes Routing zwischen ihnen erforderlich. Auf einem einzelnen Host
    übernimmt Docker diese Aufgabe mit seinen internen Netzwerken. Die
    Segmentierung lässt sich dadurch an einer Stelle definieren (siehe
    Abschnitt~\ref{sec:apphost-netze}).
  \item \textbf{Kein Vendor-Lock-in auf Proxmox.} Auf die \ac{SDN}- und
    Netzwerk-Features von Proxmox wurde bewusst verzichtet, denn Beispiele
    wie VMware zeigen, wohin eine starke Bindung an einen einzelnen
    Hersteller führen kann. Da die gesamte Netzwerksegmentierung innerhalb
    des Apphosts in Docker stattfindet, bleibt dieser portabel. Er ließe sich
    ohne Änderung der Konfiguration auch direkt auf Bare Metal oder unter
    einem anderen Hypervisor betreiben.
\end{itemize}

Die Isolation zwischen den Diensten wird folglich nicht über viele \acp{VM}
erreicht, sondern über getrennte Docker-Netzwerke, User-Namespace-Remapping
und eine strikt gehärtete Host-Konfiguration. Diese Kombination stellt den
besseren Kompromiss aus Sicherheit und Wartbarkeit dar.

\subsection{Aufbau der Entwicklungs- und Testumgebung}

Für eine realistische Entwicklung und Erprobung wurde die Umgebung des
Mandanten möglichst genau nachgebildet. Dazu wurde ein dedizierter Server (ein
AX41 bei Hetzner) angemietet und darauf Proxmox Virtual Environment
installiert. Innerhalb von Proxmox entstand ein klassisches Heimnetzwerk, wie
es bei einer Privatperson realistischerweise anzutreffen ist.

\subsubsection*{OPNsense als Router}

Als Router und Firewall des simulierten Heimnetzes diente OPNsense in einer
eigenen \ac{VM}. OPNsense ist eine quelloffene Firewall und wäre auch im
realen Einsatz beim Mandanten eine sinnvolle Wahl. Nur die OPNsense-\ac{VM}
hatte direkten Zugriff auf das Internet und maskierte den ausgehenden Verkehr
des internen Netzes per \ac{NAT}. Dahinter entsprach der Aufbau einem
gewöhnlichen Heimnetz mit \ac{DHCP} für die internen Geräte, einem lokalen
\ac{DNS}-Cache und einem flachen internen Netzwerk.

\subsubsection*{Zugang über WireGuard im Split-Tunnel}

Um so arbeiten zu können, als befänden sich die Arbeitsgeräte physisch im
Heimnetz des Mandanten, wurde auf der OPNsense ein WireGuard-Server
eingerichtet. Die Geräte des Teams verbinden sich per WireGuard im
Split-Tunnel-Modus. Nur der Verkehr in das interne Netz läuft durch den
Tunnel, während der übrige Internetverkehr der Geräte unberührt bleibt. Die
internen Dienste waren damit genauso erreichbar wie für ein Gerät im Heimnetz,
ohne dass potentiell unsichere Ports nach außen geöffnet werden mussten.

\subsubsection*{Absicherung des Proxmox-Hosts} \label{sec:SSH}

Der Proxmox-Host selbst wurde vor Beginn der eigentlichen Arbeit abgesichert.
Dazu gehören Zwei-Faktor-Authentifizierung für die Weboberfläche, komplexe
Passwörter, ausschließlich schlüsselbasierter \ac{SSH}-Zugang und die
Deaktivierung nicht benötigter Dienste. Für die Härtung wurde das Skript
\texttt{scripts/proxmox-harden.sh} entwickelt. Es deaktiviert unter anderem
die Passwort-Authentifizierung für \ac{SSH}, erzwingt moderne Cipher-Suites,
sperrt \ac{USB}-Eingaben und prüft den Secure-Boot-Status. Die Details sind in
der Installationsanleitung beschrieben.

\subsubsection*{Annahme zum Heimnetzwerk}

Die Aufgabenstellung spezifiziert das Heimnetzwerk nicht näher. Es wurde daher
ein klassisches, flaches Heimnetz angenommen und bewusst nicht weiter
verkompliziert. Der Schwerpunkt der Absicherung liegt auf dem Apphost und
seinen Diensten, also genau dort, wo die sensiblen Daten des Mandanten liegen.

\subsubsection*{Getrennte Umgebung für \ac{KI}-Experimente}

Im Rahmen des zweiten Aufgabenteils (Siehe Abschnitt~\ref{chap:ki-evaluation}-Evaluation) sollten Agenten
teilweise autonom in Proxmox arbeiten können. Damit diese Experimente die
manuelle Arbeit nicht gefährden konnten, erhielten sie eigene
Proxmox-Benutzer und vollständig getrennte Resource Pools. Bis auf einen
gemeinsam genutzten Storage-Pool waren diese Umgebungen von der eigentlichen
Infrastruktur getrennt. Die Agenten konnten dadurch nicht auf die produktive
Konfiguration zugreifen und arbeiteten ausschließlich in ihrem eigenen
Bereich.

\subsection{Der Apphost-Stack im Detail}

Der Apphost besteht aus zwei klar voneinander getrennten Schichten, dem
Betriebssystem (NixOS) und der Anwendungsschicht (Docker Compose). Beide sind
vollständig als Code beschrieben und liegen versioniert im Git-Repository.

\subsubsection{Warum NixOS}
\label{sec:warum-nixos}

Als Betriebssystem für den Apphost wurde NixOS gewählt. Diese Wahl deckt sich
in mehreren Punkten mit den Anforderungen der Aufgabe:

\begin{itemize}
  \item \textbf{Vollständig deklarativ.} Der gesamte Systemzustand, von den
    installierten Paketen über die Firewall bis zu jeder einzelnen
    Sicherheitseinstellung, ist in Konfigurationsdateien beschrieben. Genau
    das ist mit Infrastructure-as-Code gemeint, denn der Mandant kann den
    kompletten Aufbau vor der Ausführung im Repository prüfen.
  \item \textbf{Reproduzierbar.} Aus derselben Konfiguration entsteht stets
    dasselbe System, sodass sich über die Zeit kein Konfigurations-Drift
    einschleicht. Der Mandant kann die Installation auf eigener Hardware
    durchführen und erhält exakt dasselbe Ergebnis.
  \item \textbf{Atomare Updates und Rollbacks.} Jede Änderung erzeugt eine
    neue Generation. Schlägt ein Update fehl, kann beim nächsten Neustart
    eine ältere Generation aus dem Bootmenü ausgewählt werden. Das erhöht die
    Wartbarkeit deutlich.
  \item \textbf{Härtung an einer Stelle.} Alle Härtungsmaßnahmen sind zentral
    in wenigen Nix-Modulen (\texttt{security.nix}, \texttt{networking.nix},
    \texttt{docker.nix}, \texttt{secureboot.nix}) zusammengefasst. Der
    Sicherheitsstand bleibt dadurch nachvollziehbar und überprüfbar, statt
    sich über viele manuelle Schritte zu verteilen.
\end{itemize}

Die Festplatte wird über das Modul \texttt{disko} ebenfalls deklarativ
partitioniert. Zum Einsatz kommt Btrfs mit zstd-Kompression und getrennten
Subvolumes für \texttt{/}, \texttt{/nix}, \texttt{/var}, \texttt{/opt} und
\texttt{/tmp}. Das \texttt{/tmp}-Subvolume wird mit den Optionen
\texttt{nosuid}, \texttt{nodev} und \texttt{noexec} eingebunden. Der Swap
wird bei jedem Start mit einem zufälligen Schlüssel verschlüsselt, damit
keine sensiblen Daten ungeschützt im Swap verbleiben.

Optional lässt sich die Root-Partition mit \ac{LUKS}2 verschlüsseln. \label{sec:LUKS2_apphost}
Standardmäßig ist diese Option deaktiviert, da angenommen wird, dass der
gesamte Proxmox-Host vom Nutzer mit Festplattenverschlüsselung installiert
wurde. Die zusätzliche Schutzschicht kann bei der Installation aktiviert
werden. Die Umschaltung wirkt nur bei einer Neuinstallation, weil dafür die
Festplatte neu partitioniert wird.

\subsubsection{Warum Docker Compose}
Die Dienste selbst laufen als Container und werden über Docker Compose
orchestriert. Nach dem Verzicht auf Kubernetes ist das die naheliegende Wahl.
Die Konfiguration ist kompakt, gut lesbar und für den Mandanten leicht
nachvollziehbar. Der gesamte Stack wird über eine zentrale
\texttt{docker-compose.yml} zusammengesetzt, die per \texttt{include}
einzelne, thematisch gruppierte Compose-Dateien einbindet (zum Beispiel
\texttt{media}, \texttt{documents}, \texttt{monitoring}, \texttt{cloud}).
Jede Datei bleibt dadurch überschaubar und einem klaren Zweck zugeordnet.

\subsubsection{Netzwerksegmentierung}
\label{sec:apphost-netze}
Statt alle Container in ein gemeinsames Netzwerk zu legen, erhält jede
Dienstgruppe ihr eigenes Docker-Netzwerk. Eine Datenbank ist nur in dem Netz
erreichbar, in dem auch die zugehörige Anwendung liegt. Diese Aufteilung
begrenzt die laterale Bewegung eines Angreifers. Selbst nach der
Kompromittierung eines einzelnen Dienstes ist von dort kein direkter Zugriff
auf alle anderen Dienste möglich.


\subsubsection{Reverse Proxy mit Traefik}
Traefik ist der einzige Einstiegspunkt für Web-Verkehr. Er terminiert
\ac{TLS} und verteilt die Anfragen anhand der Subdomain an den jeweiligen
Dienst. Die Zuordnung erfolgt über Container-Labels, die Traefik direkt aus
Docker liest. Damit Traefik keinen vollen Zugriff auf den Docker-Socket
benötigt, wird die Docker-\ac{API} ausschließlich über einen eingeschränkten
Socket-Proxy ausgelesen. Dieser lässt nur lesende Anfragen zu, sodass selbst
eine Kompromittierung von Traefik keinen Vollzugriff auf das System
ermöglicht.

Für den internen Zugriff über das Heimnetz wird eine eigene Domain genutzt (in
den Tests \path{lbaecker.de}), deren \ac{DNS}-A-Eintrag direkt auf die private
\ac{IP}-Adresse des Servers im Heimnetzwerk zeigt. Von außen ist dieser
Eintrag zwar öffentlich auflösbar, jedoch nicht erreichbar, da die private
\ac{IP} außerhalb des Heimnetzes nicht routbar ist. Angreifer können auf
diesem Weg lediglich die interne \ac{IP}-Adresse des Apphosts ermitteln.
Diese geringfügige Informationspreisgabe stellt einen bewussten Tradeoff dar.
Alternativ ließe sich ein statischer \ac{DNS}-Eintrag im lokalen Resolver des
Routers hinterlegen. Ein sicherheitsbewusster Nutzer wie der Mandant besitzt
jedoch voraussichtlich Geräte, die standardmäßig auf sicheres \ac{DNS} setzen
und den via \ac{DHCP} oder Router Advertisement verteilten \ac{DNS}-Server
ignorieren. Mit dem gewählten Ansatz lässt sich die \ac{IP}-Adresse mit jedem
\ac{DNS}-Server auflösen. Ein zusätzlicher interner Eintrag im Router bleibt
dennoch möglich.

Der Vorteil einer echten Domain liegt darin, dass ein öffentlich
vertrauenswürdiges Zertifikat von Let's Encrypt bezogen werden kann und keine
selbstsignierten Zertifikate erforderlich sind. Daraus ergeben sich mehrere
praktische Vorteile. Zunächst treten keine Zertifikatswarnungen im Browser
auf. Das ist auch sicherheitsrelevant, denn wer sich angewöhnt,
Zertifikatswarnungen wegzuklicken, übersieht irgendwann auch eine echte
Man-in-the-Middle-Warnung. Ein gültiges Zertifikat erhält diese Warnung als
verlässliches Signal. Außerdem funktioniert die volle Vertrauenskette ohne
manuelles Verteilen einer eigenen \ac{CA} auf jedes Gerät des Mandanten.
Gerade bei Mobiltelefonen und bei Apps, die selbstsignierte Zertifikate
strikt ablehnen, verursacht das sonst häufig Probleme. Schließlich ist die
Verwaltung vollständig automatisiert, da Traefik die Zertifikate rechtzeitig
selbst erneuert und somit keine abgelaufenen Zertifikate entstehen.

Die Zertifikate bezieht Traefik über \texttt{Let's Encrypt} mit der
\ac{DNS}-01-Challenge. Der große Vorteil dieser Methode besteht darin, dass
der Server dafür nicht aus dem Internet erreichbar sein muss. Über die
\ac{DNS}-\ac{API} des Anbieters (hier \texttt{Cloudflare}) wird lediglich kurzzeitig
ein TXT-Eintrag gesetzt. Die Ports 80 und 443 können vollständig hinter dem
Router verbleiben. An \texttt{Cloudflare} fließen dabei keine Inhalte, sondern nur der
Nachweis der Kontrolle über die Domain.

Eine bewusste Ausnahme bildet der Tor-Zugang. Onion-Adressen lassen sich von
Let's Encrypt nicht zertifizieren. Da die Verbindung über Tor bereits
Ende-zu-Ende verschlüsselt ist, wird dort ein selbstsigniertes Zertifikat
verwendet, was im Tor-Browser unkritisch ist. Auch wenn ein Zertifikat nicht
zwingend notwendig wäre, bildet es eine zusätzliche Schutzschicht für den
Fall, dass das Tor-Netzwerk kompromittiert wird oder Schwächen im
Tor-Protokoll auftreten.

Die zentrale Zertifikatsverwaltung kommt auch \ac{MQTT} zugute (hier
Mosquitto). Der Broker wird nicht über einen eigenen \ac{TLS}-Listener nach
außen exponiert, sondern über einen zusätzlichen \ac{TCP}-Entrypoint von
Traefik getunnelt. Traefik terminiert das \ac{TLS} anhand des im
\ac{TLS}-ClientHello mitgesendeten Hostnamens (\ac{SNI}) mit demselben
Let's-Encrypt-Zertifikat, das auch für den regulären Web-Zugriff verwendet
wird. Die entschlüsselten Daten werden innerhalb des Docker-Netzes an den
internen, unverschlüsselten \ac{MQTT}-Listener von Mosquitto weitergereicht.
Auf dem Apphost muss damit kein separates Zertifikat für Mosquitto erzeugt,
verteilt und erneuert werden. Alle Geräte des Mandanten vertrauen diesem
öffentlich signierten Zertifikat in der Regel von Haus aus.

\subsubsection{Single Sign-On mit Authelia}
\label{sec:apphost-authelia}
Authelia ist der zentrale Authentifizierungsdienst und schützt die Dienste,
die keine ausreichend starke eigene Anmeldung mitbringen. Je nach
Unterstützung des jeweiligen Dienstes kommen zwei Mechanismen zum Einsatz.

\begin{itemize}
  \item \textbf{\ac{OIDC}} für Dienste mit nativer Unterstützung, zum
    Beispiel Grafana, Immich, Paperless und Forgejo. Diese Dienste leiten zur
    Anmeldung an Authelia weiter und vertrauen anschließend dem ausgestellten
    Token. Die Benutzeridentität wird dadurch sauber in den Dienst
    übernommen.
  \item \textbf{Forward-Auth} für Dienste ohne brauchbare eigene
    Zugriffskontrolle, zum Beispiel Prometheus, Alertmanager und die
    Homepage. Traefik prüft dabei bei jeder Anfrage zuerst bei Authelia, ob
    der Nutzer angemeldet ist, bevor die Anfrage den Dienst erreicht.
\end{itemize}

Die Benutzer und ihre Passwörter liegen in einer Datei mit Argon2id-Hashes,
die beim Setup automatisch generiert wird. Die zentrale Anmeldung ist damit
deklarativ beschrieben und nicht in einer von außen nicht prüfbaren Datenbank
abgelegt.

\subsubsection{Zensurresistenter Zugriff über Tor} \label{sec:Tor}

Auf ausdrücklichen Kundenwunsch ist ein Großteil der Dienste zusätzlich über
das Tor-Netzwerk als Onion Service v3 erreichbar. Dafür läuft ein
Tor-Container, der einen Onion Service bereitstellt und den Verkehr an einen
eigenen Traefik-Einstiegspunkt weiterreicht. Der Mandant kann dadurch auch
aus stark überwachten oder zensierten Netzen auf seine Dienste zugreifen.
Metadaten fließen dabei nicht an den lokalen Provider ab und es muss kein
Port nach außen geöffnet werden. Der Tor-Container ist als reiner Client
konfiguriert (kein Relay, kein Exit, deaktivierter Control- und
\ac{SOCKS}-Port).

Praktisch entsteht damit ein dualer \ac{URL}-Raum. Im Heimnetz sind die
Dienste unter \texttt{*.domain.tld} mit gültigem Zertifikat erreichbar. Über
Tor sind sie unter \texttt{*.<onion-service>.onion} erreichbar, wobei das
oben beschriebene zentrale selbstsignierte Wildcard-Zertifikat genutzt wird.

Die Tor-Erreichbarkeit lässt sich \emph{pro Dienst} explizit ein- und
ausschalten. Jeder Dienst erhält neben seinen normalen Traefik-Labels ein
zusätzliches Set an Labels für den Tor-Einstiegspunkt. Ohne dieses Set ist
der Dienst nicht über Tor erreichbar, während sich am internen Zugriff nichts
ändert. Der Mandant kann somit gezielt steuern, welche Dienste anonym von
außen verfügbar sind und welche bewusst im Heimnetz verbleiben. Collabora ist
aus Sicherheitsgründen standardmäßig nicht über Tor freigegeben (siehe
Abschnitt~\ref{sec:kompromisse_apphost}). Auch Garage hat bewusst keinen
Tor-Einstiegspunkt, denn als \acs{S3}-kompatibler Objektspeicher im
Hintergrund der anderen Dienste hat es keine eigenen Endnutzer, die von
außerhalb des Heimnetzes darauf zugreifen müssten.

\paragraph{Kompatibilität der Anwendungen und \ac{OIDC}.} \mbox{}\\
Nicht jede Anwendung kommt mit zwei verschiedenen \acp{URL} gleich gut
zurecht. Manche verhalten sich völlig transparent. Home Assistant etwa
orientiert sich am jeweils anfragenden Host und stellt sich automatisch auf
die gerade verwendete \ac{URL} ein, sodass der Wechsel zwischen internem Pfad
und Tor dort unbemerkt funktioniert. Eine echte Einschränkung betrifft
dagegen die Dienste, die sich über \ac{OIDC} an Authelia anmelden. \ac{OIDC}
bindet die ausgestellten Tokens, die Redirect-\acp{URI} und vor allem die
sogenannte Issuer-\ac{URL} fest an eine einzige Basis-Adresse. Wird Authelia
intern unter \texttt{auth.domain} betrieben, sind alle Tokens auf genau diese
Adresse ausgestellt. Meldet sich derselbe Nutzer anschließend über die
Tor-Adresse an, schlägt die Validierung fehl. Die Issuer-\ac{URL} in der
Antwort passt dann nicht mehr zu der Adresse, über die der Zugriff erfolgt.
Ein unmittelbarer Wechsel zwischen internem Pfad und Tor-Pfad ist bei
\ac{OIDC}-geschützten Diensten daher nicht möglich. Bei reinen
Forward-Auth-Diensten (zum Beispiel Prometheus oder die Homepage) tritt das
Problem nicht auf, weil dort kein Token an eine feste Issuer-\ac{URL}
gebunden wird. Es handelt sich um eine reine Komforteinschränkung und nicht
um ein Sicherheitsproblem. Der Umgang damit und die ergänzende Empfehlung
eines WireGuard-Zugangs werden in Abschnitt~\ref{sec:kompromisse_apphost}
erläutert.

Der Zusammenhang mit dem \ac{HTTP}-Host-Header ist direkt. Authelia bestimmt
die Issuer-\ac{URL} aus einer statisch konfigurierten Basis-Adresse, prüft
aber bei jeder Abfrage, ob der Host- beziehungsweise
X-Forwarded-Host-Header zu dieser passt. Greift ein Client über eine andere
Adresse zu, etwa die Tor-Adresse, schlägt dieser Abgleich fehl. Selbst wenn
der Abgleich umgangen werden könnte, bliebe das Problem bestehen. Der
ausgestellte \ac{ID}-Token trüge weiterhin die falsche Issuer-\ac{URL} und
die Anwendung validiert sie gegen die hinterlegte Adresse. Eine Lösung über
den Browser, etwa durch Header-Manipulation, ist daher nicht sinnvoll
möglich. Würden die \ac{OIDC}-Klienten Host-Header nicht strikt validieren,
wären sie für Injection-Angriffe anfällig. Die feste Bindung ist folglich
kein Mangel, sondern eine bewusste Sicherheitsfunktion von \ac{OIDC}. Sie
lässt sich nicht umgehen, ohne die Sicherheit des Gesamtsystems erheblich
herabzusetzen.

\subsubsection{Monitoring und Alerting} \label{sec:apphost-monitoring}

Für die Überwachung kommt der etablierte Stack aus Prometheus (Metriken),
Loki (Logs), Grafana (Dashboards) und Alertmanager (Alarmierung) zum Einsatz.
Node-Exporter und cAdvisor liefern Host- und Container-Metriken. Grafana
Alloy sammelt die Logs ein, sowohl die Docker-Container-Logs als auch
\texttt{syslog} und \texttt{auth.log} des Hosts, und schreibt sie nach Loki.
Prometheus und Loki arbeiten rein intern. Nur Grafana, Prometheus und
Alertmanager sind über den Reverse Proxy zugänglich und dort durch Authelia
geschützt.

Dieser Stack stammt ursprünglich aus der Kubernetes-Phase, denn im
cloud-nativen Umfeld ist die Kombination aus Prometheus und Grafana der
De-facto-Standard für Monitoring. Beim Umbau auf Docker Compose konnte er
ohne größere Anpassungen übernommen werden. Genau darin liegt einer seiner
größten Vorzüge, denn die Komponenten sind nicht an Kubernetes gebunden,
sondern laufen ebenso als gewöhnliche Container. Das einmal erarbeitete
Monitoring-Konzept musste daher nicht verworfen werden:

\begin{itemize}
  \item \textbf{Pull-basiert und einfach.} Prometheus holt die Metriken
    selbstständig von den Exportern ab. Es muss nichts an einen externen
    Dienst übertragen werden und alle Daten bleiben auf dem Host.
  \item \textbf{Geringer Ressourcenbedarf.} Gerade Loki ist bewusst
    leichtgewichtig gehalten. Es indexiert nur Metadaten statt des vollen
    Textes wie etwa Elastic oder OpenSearch und passt damit gut zu einem
    einzelnen, sparsam ausgelegten Host.
  \item \textbf{Alles aus einer Oberfläche.} Grafana führt Metriken und Logs
    in einer einzigen Oberfläche zusammen, sodass der Mandant nicht mehrere
    Werkzeuge bedienen muss.
  \item \textbf{Flexible Alarmierung.} Der Alertmanager trennt die Auswertung
    der Alarme von deren Zustellung und lässt sich frei konfigurieren, in
    diesem Fall in Richtung des selbst gehosteten ntfy.
\end{itemize}

Besonders wichtig für das Bedrohungsmodell ist die Alarmierung. Die Aufgabe
verlangt ausdrücklich, dass Push-Benachrichtigungen nicht über fremde
Cloud-Dienste laufen, die Metadaten abgreifen könnten. Deshalb kommt das
selbst gehostete \texttt{ntfy} zum Einsatz. Der Alertmanager schickt seine
Alarme per Webhook an \texttt{ntfy}. Der Mandant empfängt sie sowohl am Desktop als
auch auf seinem \texttt{GrapheneOS}-Smartphone. Da \texttt{ntfy} selbst gehostet ist, verlassen
die Benachrichtigungen die eigene Infrastruktur nicht. Definierte Alarme
betreffen unter anderem hohe \ac{CPU}-Last, knappen Speicherplatz,
ausgefallene Container und ein nicht erreichbares Backend hinter \texttt{Traefik}.

\subsubsection{Absicherung der Container}
\label{sec:apphost-container}

Die Netzwerksegmentierung allein reicht nicht aus. Damit ein einzelner
kompromittierter Dienst möglichst wenig Schaden anrichten kann, wird jeder
Container nach demselben Grundmuster gehärtet. Die Einstellungen stehen
direkt in den Compose-Dateien und sind damit für den Mandanten
nachvollziehbar:

\begin{itemize}
  \item \textbf{Keine neuen Privilegien.} Jeder Container läuft mit\\
    \texttt{no-new-privileges:true}. Ein Prozess im Container kann seine
    Rechte damit nicht nachträglich erhöhen, etwa über setuid-Binaries.
  \item \textbf{Capabilities auf null und gezielt zurück.} Mit
    \texttt{cap\_drop: ALL} werden zunächst sämtliche Linux-Capabilities
    entfernt. Anschließend werden pro Dienst nur die wenigen wieder
    freigegeben, die er tatsächlich benötigt. Typischerweise sind das SETUID
    oder SETGID, wenn der Container intern den Benutzer wechselt.
  \item \textbf{Read-only Dateisystem.} Wo der Dienst es zulässt (etwa
    Traefik, Tor, Garage, Node-Exporter, Homepage, ntfy und Mosquitto), wird
    das Wurzeldateisystem des Containers schreibgeschützt eingebunden. Für
    die wenigen Stellen mit echtem Schreibbedarf dient ein kleines, nicht
    ausführbares tmpfs im \ac{RAM}.
  \item \textbf{Ressourcenlimits.} Jeder Container besitzt \ac{CPU}- und
    Speicherobergrenzen. Ein einzelner, möglicherweise kompromittierter
    Dienst kann den Host dadurch nicht lahmlegen.
  \item \textbf{Seccomp und AppArmor.} Es gilt das Standard-Seccomp-Profil
    von Docker, das gefährliche Syscalls blockiert. Auf Host-Ebene erzwingt
    AppArmor zusätzliche Mandatory-Access-Control-Profile.
  \item \textbf{User-Namespace-Remapping.} Auf Daemon-Ebene wird root im
    Container auf einen unprivilegierten Bereich auf dem Host abgebildet.
    Selbst wer im Container root wird, bleibt auf dem Host ein Benutzer ohne
    Rechte.
\end{itemize}

Nicht jeder Container folgt diesem Muster vollständig. Neben Collabora, das für
das Rendern von Dokumenten erweiterte Rechte benötigt und im Abschnitt~\ref{sec:kompromisse_apphost} 
behandelt wird, brauchen auch einige Monitoring-Komponenten und
die Docker-Anbindung selbst gezielte Ausnahmen:

\begin{itemize}
  \item \textbf{cAdvisor} läuft mit \texttt{privileged: true} und
    \texttt{userns\_mode: host}, weil es zur Erhebung von Container-Metriken
    lesend auf cgroups, den Kernel-Ringpuffer (\texttt{/dev/kmsg}) und
    interne Docker-Zustände zugreifen muss. Ohne erweiterte Rechte sind diese
    nicht sichtbar.
  \item \textbf{Node-Exporter} läuft mit Host-PID-, Host-Netzwerk- und
    Host-User-Namespace (\texttt{pid: host}, \texttt{network\_mode: host},
    \texttt{userns\_mode: host}), weil er Prozesse und
    Netzwerkschnittstellen des Hosts selbst vermisst und nicht die eines
    isolierten Containers.
  \item \textbf{socket-proxy} läuft mit \texttt{userns\_mode: host} und
    bindet den Docker-Socket read-only ein. Der Socket funktioniert
    ausschließlich mit den Host-\acp{UID}, mit denen der Docker-Daemon selbst
    läuft. Unter User-Namespace-Remapping wären die gemappten \acp{UID} im
    Container nicht identisch mit denen, die der Socket erwartet. Traefik
    selbst erhält im Gegenzug keinen direkten Socket-Zugriff und spricht
    ausschließlich den lesebeschränkten Socket-Proxy über das Netzwerk an
    (siehe oben).
  \item \textbf{Grafana Alloy} bindet den Docker-Socket ebenfalls read-only
    ein, um Container-Logs einsammeln zu können. Es läuft dabei ohne
    \texttt{userns\_mode: host} und mit vollständig entfernten Capabilities.
\end{itemize}

Alle diese Ausnahmen sind eng auf das jeweils benötigte Recht begrenzt und
betreffen bewusst nur Monitoring- und Infrastrukturkomponenten. Diese
verarbeiten keine Nutzdaten des Mandanten, sondern beobachten Host-
beziehungsweise Container-Zustand oder leiten Anfragen weiter. Die
produktiven Dienste des Mandanten bleiben vom vollen Härtungsmuster erfasst.

Ergänzt wird die Container-Härtung durch automatische Sicherheitsscans. Ein
geplanter Job prüft die laufenden Images wöchentlich (montags um 02:00 Uhr)
mit Trivy auf bekannte Schwachstellen der Stufen HIGH und CRITICAL und legt
den Bericht unter \textit{/var/log/docker-security-scan.log} ab. Zusammen mit
den durch RenovateBot laufend aktuell gehaltenen Images entsteht ein
kontinuierlicher Kreislauf. Neue Schwachstellen werden zeitnah erkannt und
das passende Update steht meist bereits als Pull Request bereit.


\subsubsection{Härtung von Webanfragen mit Traefik}
Weil Traefik der einzige Einstiegspunkt für Web-Verkehr ist, werden dort auch
die Härtungsmaßnahmen gebündelt, die für den \ac{HTTPS}-Verkehr aller Dienste
gelten. Statt Security-Header, \ac{TLS}-Einstellungen und Zugriffsschutz in
mehr als zwanzig einzelnen Diensten separat zu pflegen, werden sie zentral
als Traefik-Middleware definiert und über sogenannte Chains an die jeweiligen
Router gebunden. Ein neuer Dienst erbt damit automatisch den vollen Schutz,
sobald er die passende Kette an Middlewares referenziert.

\begin{itemize}
  \item \textbf{\ac{TLS}-Konfiguration nach \ac{BSI} \ac{TR}-02102-2.}
    Aktuell gilt für sämtliche Entrypoints ausschließlich das
    \texttt{default}-\ac{TLS}-Profil mit \ac{TLS}~1.3. Daneben liegt ein
    zweites, ebenfalls konformes Profil (\texttt{intermediate}) vor. Es
    erlaubt \ac{TLS}~1.2 nur mit den vom \ac{BSI} empfohlenen
    \ac{AEAD}-Cipher-Suites (AES-GCM, ChaCha20-Poly1305) sowie den Kurven
    X25519, P-256 und P-384 als Fallback. Da kein Router explizit auf
    \texttt{tls.options=intermediate} verweist, bleibt dieses Fallback-Profil
    derzeit ungenutzt und alle Clients müssen \ac{TLS}~1.3 sprechen. Bei
    Bedarf lässt es sich für einzelne Router aktivieren, etwa falls ein
    Altgerät des Mandanten kein \ac{TLS}~1.3 unterstützt. Die Option
    \texttt{sniStrict} sorgt in beiden Profilen dafür, dass Anfragen ohne
    passenden Hostnamen abgewiesen werden, statt auf ein Default-Zertifikat
    auszuweichen.
  \item \textbf{Erzwungenes \ac{HTTPS} und \ac{HSTS}.} Der
    \ac{HTTP}-Entrypoint leitet jede Anfrage permanent auf \ac{HTTPS} um.
    Zusätzlich setzt Traefik einen \ac{HSTS}-Header mit zwei Jahren
    Gültigkeit und Preload-Flag für alle Subdomains, damit Browser gar nicht
    erst versuchen, sich unverschlüsselt zu verbinden. Browser sind in den
    letzten Jahren zwar auf ein \ac{HTTPS}-by-default-Modell gewechselt,
    \ac{HSTS} erzwingt \ac{HTTPS} jedoch dauerhaft. Das Eintragen in die
    Preload-Listen der Browser erfordert einen öffentlich erreichbaren
    Webserver und einen Antrag auf \texttt{hstspreload.org}. Das
    Preload-Flag hat daher momentan keine Wirkung und stellt eine
    Vorbereitung dar.
  \item \textbf{Security-Header für alle Dienste.} Über die Middleware\\
    \texttt{security-headers} setzt Traefik sicherheitsrelevante Header wie
    \texttt{Content-Security-Policy}, \texttt{X-Frame-Options},\\
    \texttt{X-Content-Type-Options: nosniff}, eine restriktive Referrer- und
    Permissions-Policy (deaktiviert unter anderem Kamera, Mikrofon und
    Standort) sowie Cross-Origin-Opener-, Embedder- und Resource-Policy.
    Gleichzeitig werden \texttt{Server}- und \texttt{X-Powered-By}-Header
    geleert, damit Angreifer aus einer Antwort nicht ableiten können, welche
    Software in welcher Version dahintersteht. Ein
    \texttt{X-Robots-Tag: noindex} verringert außerdem das Risiko, dass
    interne Dienste in Suchmaschinen erscheinen.
  \item \textbf{Schutz vor gefälschten Auth-Headern.} Dienste wie Grafana
    oder Paperless vertrauen im Onion-Pfad einem
    \texttt{Remote-User}-Header, den Authelia nach erfolgreicher Prüfung
    setzt. Über das Clearnet darf dieser Header niemals vom Client selbst
    stammen. Die Middleware \texttt{strip-remote-user} entfernt deshalb auf
    allen normalen Routern jeden vom Client mitgeschickten
    \texttt{Remote-*}-Header, bevor die Anfrage den Dienst erreicht. Niemand
    kann sich somit durch einen gefälschten Header als anderer Benutzer
    ausgeben.
  \item \textbf{Auswertbare Logs.} Das Access-Log läuft im \ac{JSON}-Format
    und wird an Loki weitergereicht. Sensible Header wie
    \texttt{Authorization} und \texttt{Cookie} werden dabei grundsätzlich
    verworfen. \texttt{User-Agent} und \texttt{X-Forwarded-For} bleiben für
    mögliche zukünftige forensische Auswertungen erhalten.
\end{itemize}

Diese Maßnahmen ergänzen die oben beschriebene Container-Härtung von Traefik
selbst, unter anderem \texttt{cap\_drop: ALL}, ein read-only Dateisystem und
den Zugriff auf die Docker-\ac{API} ausschließlich über den Socket-Proxy.


\subsubsection{Umgang mit Geheimnissen}
Alle Geheimnisse (Passwörter, Token, Schlüssel) liegen in einer
\texttt{.env}-Datei und einem \texttt{secrets/}-Verzeichnis. Beide werden
nicht in Git eingecheckt. Aus den Klartextwerten der \texttt{.env} erzeugen
kleine Skripte die tatsächlich benötigten gehashten Werte, etwa die
Argon2id-Hashes für Authelia, die bcrypt-Hashes für ntfy, die Passwortdatei
für Mosquitto sowie die \texttt{rpcauth}-Hashes für den Bitcoin-Core-Node des
Hot-Wallet-Stacks. Das Verzeichnis \texttt{secrets/} ist auf
Dateisystemebene strikt geschützt und nur für den \texttt{apphost}-Benutzer
zugänglich. Ändert der Mandant ein Passwort, generiert er die Hashes mit
einem einzigen Befehl neu.

\subsubsection{Kontrollierte Updates mit Renovate}
\label{sec:renovate}
Ein wiederkehrendes Problem beim Self-Hosting ist, dass Container-Images
veralten und so unbemerkt Sicherheitslücken offen bleiben. Um das zu
verhindern, überwacht RenovateBot das Repository und öffnet automatisch einen
Pull Request, sobald für ein Image eine neue Version verfügbar ist. Daraus
ergibt sich ein kontrollierter und nachvollziehbarer Update-Prozess.

\begin{itemize}
  \item Jedes Update kommt als einzelner, dokumentierter Pull Request und
    geschieht nicht unbemerkt im Hintergrund.
  \item Unkritische Updates für den Monitoring-Stack werden automatisch
    zusammengefasst und können automatisch gemerged werden.
  \item Sicherheitskritische Komponenten wie Datenbanken oder Traefik werden
    bewusst nicht automatisch gemerged, sondern erst nach einer Wartezeit und
    manueller Prüfung übernommen.
  \item Major-Updates mit möglichen Breaking Changes warten länger und werden
    stets von Hand geprüft.
\end{itemize}

Die genannte Wartezeit wird über die Renovate-Option
\textit{minimumReleaseAge} umgesetzt. Für Major-Updates beträgt sie 14~Tage
und für das exponierte Traefik-Image 7~Tage. Renovate öffnet den
Update-Pull-Request zwar sofort, markiert den zugehörigen Status-Check aber
als ausstehend, bis die konfigurierte Zeit seit dem Release verstrichen ist.
Der Sicherheitsgewinn dieser Verzögerung ist beträchtlich, denn wer ein
Release sofort übernimmt, bringt ungeprüften Code direkt in Produktion.
Kompromittierte Releases werden in der Praxis meist innerhalb von Stunden bis
wenigen Tagen von der Community entdeckt und vom jeweiligen
Registry-Betreiber zurückgezogen. Eine kurze Wartezeit fängt sie daher in
aller Regel ab. Ein anschauliches Beispiel lieferte der Vorfall um das
NPM-Paket \texttt{ua-parser-js} im Oktober 2021. Nach der Übernahme des
Maintainer-Accounts wurden drei bösartige Versionen veröffentlicht, die auf
Windows- und Linux-Systemen Zugangsdaten stahlen und einen Krypto-Miner
nachluden \cite{cisa_ua_parser_js_2021}.

Nach dem Mergen eines Pull Requests genügt auf dem Server ein einziger
Befehl, um die neuen Images zu ziehen und die Container neu zu starten.
Ergänzt wird das durch einen wöchentlichen Schwachstellen-Scan der laufenden
Images mit Trivy und die tägliche automatische Aktualisierung des
NixOS-Systems. Die Angriffsfläche bleibt so dauerhaft klein, ohne dass
manuell nach neuen Versionen gesucht werden muss.


\subsubsection{Hash-Pinning der Container-Images}
\label{sec:apphost-digest-pinning}

Ein reiner Versions-Tag wie \texttt{traefik:v3.7.1} ist kein unveränderlicher
Bezeichner. Er kann nachträglich auf einen anderen Image-Inhalt umgebogen
werden, sei es versehentlich oder durch Kompromittierung. Für den
Reproduzierbarkeitsanspruch, der bereits für NixOS formuliert wurde
(siehe Abschnitt~\ref{sec:warum-nixos}), stellt das einen Bruch dar. Dieselbe
\texttt{docker-compose.yml} könnte zu unterschiedlichen Zeitpunkten
unterschiedliche Images ziehen, ohne dass sich am Tag etwas ändert.

Die Image-Referenzen werden deshalb zusätzlich zum Tag um den SHA-256-Digest
des jeweiligen Manifests ergänzt, also zum Beispiel
\texttt{traefik:v3.7.1@sha256:<digest>}. Der Digest ist inhaltsadressiert und
zeigt exakt auf einen einzigen, unveränderlichen Layer-Stand. Ein
\texttt{docker compose pull} liefert damit garantiert dasselbe Image, das zum
Zeitpunkt des jeweiligen Commits vorlag, unabhängig davon, worauf der Tag im
Upstream-Repository später zeigt.

Damit dieses Pinning keinen zusätzlichen manuellen Pflegeaufwand verursacht,
übernimmt RenovateBot (siehe Abschnitt~\ref{sec:renovate}) auch diesen Teil.
Renovate ergänzt für jedes Image automatisch den passenden Digest. Bei jeder
neuen Version und jedem neuen Digest desselben Tags, etwa nach einem erneuten
Build desselben Upstream-Releases, öffnet es einen entsprechenden Pull
Request. Reproduzierbarkeit und automatisch gepflegte Updates schließen sich
damit nicht aus, sondern ergänzen einander.


\subsubsection{Übersicht der Dienste}

Insgesamt stellt der Apphost deutlich mehr als die geforderten 15 Dienste
bereit. Das Spektrum reicht von Medien (Jellyfin, Immich) über Dokumente
(Paperless, Bento-\ac{PDF}), Cloud und Office (OpenCloud, Collabora),
Hausautomation (Home Assistant, Mosquitto), Kommunikation (ntfy, Bichon) und
Entwicklung (Forgejo) bis zu Werkzeugen wie dem Passwortmanager Vaultwarden
und einer Startseite (Homepage). Zusätzlich läuft auf dem Apphost der
Hot-Wallet-Stack für den automatisierten Bitcoin-Zahlungsverkehr, der in
Abschnitt~\ref{Krypto} gesondert beschrieben wird. Die vollständige Liste mit
Subdomains und Beschreibungen findet sich im Anhang.

\subsection{Abdeckung des Bedrohungsmodells} \label{sec:gefahrenmodell}

Die Aufgabenstellung gibt drei zentrale Angriffsvektoren vor. Im Folgenden
wird für jeden davon gezeigt, mit welchen Maßnahmen er adressiert wird. Eine
ausführliche Auflistung aller einzelnen Einstellungen findet sich im Anhang.

\subsubsection{Remote-Angriffe}

Hierunter fallen Phishing, Brute-Force, kompromittierte Endgeräte und
Supply-Chain-Angriffe auf Abhängigkeiten.

\begin{itemize}
  \item \textbf{Keine offenen Ports.} Der Web-Zugriff erfolgt nur über das
    Heimnetz und über Tor und die Zertifikate werden per \ac{DNS}-01 bezogen.
    Es muss daher kein einziger Port ins Internet geöffnet werden, was
    klassischen Remote-Angriffen von vornherein die Angriffsfläche nimmt.
    Auch intern werden zusätzlich lediglich Ports für \ac{SSH}, \ac{HTTP}, \ac{HTTPS} und \ac{MQTT} benötigt.
  \item \textbf{Starke, zentrale Authentifizierung.} Authelia mit Argon2id
    und Multi-Faktor-Authentifizierung schützt die Dienste und bringt einen
    eigenen Brute-Force-Schutz auf Anwendungsebene mit (siehe
    Abschnitt~\ref{sec:apphost-authelia}). Ergänzend sperrt Fail2ban nach
    wenigen Fehlversuchen sowohl bei \ac{SSH} als auch bei fehlgeschlagenen
    Web-Logins auf Netzwerkebene und verlängert die Sperrzeit progressiv.
    Die Firewall blockiert eingehenden Verkehr standardmäßig und erlaubt nur
    das Nötigste. \ac{SSH} ist zusätzlich rate-limitiert.
  \item \textbf{Supply-Chain.} RenovateBot, der wöchentliche Trivy-Scan und
    die täglichen NixOS-Updates sorgen dafür, dass bekannte Schwachstellen
    schnell geschlossen werden. Wird dennoch ein Container kompromittiert,
    greift die in Abschnitt~\ref{sec:apphost-container} beschriebene
    Container-Härtung. Entfernte Capabilities, \texttt{no-new-privileges},
    read-only Dateisysteme, das Seccomp-Profil, AppArmor und vor allem das
    User-Namespace-Remapping begrenzen den möglichen Schaden, denn root im
    Container ist nicht root auf dem Host. Die Netzwerksegmentierung
    (siehe Abschnitt~\ref{sec:apphost-netze}) verhindert zusätzlich, dass sich ein
    Angreifer von einem Dienst aus frei weiterbewegen kann.
\end{itemize}

\subsubsection{Lokale Angriffe (Bad\ac{USB}, Einklinken ins Netz)}

Die Aufgabe stellt klar, dass das Aufschrauben des Servers out-of-scope ist.
Das Anstecken bösartiger \ac{USB}-Geräte und das unautorisierte Einklinken
ins Netzwerk müssen dagegen abgewehrt werden.

\begin{itemize}
  \item \textbf{\ac{USB}-Angriffe.} Die Kernel-Module für \ac{USB}-Speicher
    (\texttt{usb-storage}, \texttt{uas}) und weitere nicht benötigte
    Hardware-Module sind per Blacklist deaktiviert. Auch auf dem
    Proxmox-Host werden \ac{USB}-Eingaben gesperrt. Ein angestecktes
    Bad\ac{USB}-Gerät findet damit keinen direkten Angriffsweg.
  \item \textbf{Verifizierter Boot.} Secure Boot über lanzaboote stellt
    sicher, dass nur signierte Boot-Komponenten starten. Zusammen mit
    erzwungener Kernel-Modul-Signatur (\texttt{module.sig\_enforce}) und dem
    Kernel-Lockdown wird verhindert, dass ein lokaler Angreifer manipulierten
    Code einschleust.
  \item \textbf{\ac{DMA}-Schutz.} Die \ac{IOMMU} wird erzwungen und früher
    \ac{PCI}-\ac{DMA}-Zugriff blockiert. Angriffe über direkten
    Speicherzugriff werden dadurch erschwert.
  \item \textbf{Netzwerk.} Wer sich unautorisiert ins Netz einklinkt, trifft
    auf eine Default-Deny-Firewall und auf Dienste, die hinter Authelia
    liegen. Ein bloßer Netzwerkzugang genügt somit nicht, um an die Dienste
    zu gelangen.
\end{itemize}

\subsubsection{Metadaten-Leaks}

Damit keine sensiblen Nutzungs- oder Standortdaten an den Provider abfließen,
setzt das System an mehreren Stellen an:

\begin{itemize}
  \item \textbf{Tor.} Der Zugriff über Onion Services verbirgt sowohl die
    Identität des Servers als auch die des Nutzers vor dem lokalen Provider.
  \item \textbf{Verschlüsseltes \ac{DNS}.} Der Host nutzt
    \ac{DNS}-over-\ac{TLS} mit \ac{DNSSEC}-Validierung gegen Resolver ohne
    Logging-Policy (Cloudflare und Quad9). \ac{DNS}-Anfragen landen dadurch
    nicht im Klartext beim Provider.
  \item \textbf{Authentifizierte Zeit.} Die Zeitsynchronisation läuft über
    \ac{NTS} und damit authentifiziert über \ac{TLS}.
  \item \textbf{Keine fremde Cloud für Benachrichtigungen.}
    Push-Benachrich"=tigungen laufen, wie beschrieben, ausschließlich über das
    selbst gehostete ntfy.
\end{itemize}

\subsubsection{Eingegangene Kompromisse}\label{sec:kompromisse_apphost}

Kein reales Setup ist frei von Kompromissen. Deren offene Benennung gehört zum
Entwicklungsprozess und ist für diese Ausarbeitung von hoher Relevanz. Die
folgenden Punkte behandeln daher die bewusst eingegangenen Kompromisse.

\begin{itemize}
  \item \textbf{Ein Host statt vieler \acp{VM}.} Der Aufbau bündelt viele
    Dienste auf einem Host. Dieser ist damit ein attraktiveres Ziel, als es
    eine einzelne Mini-\ac{VM} wäre. Das erscheint vertretbar, weil die
    Isolation über Docker-Netzwerke und User-Namespaces sowie eine sehr
    umfangreiche Host-Härtung erfolgt. Der Gewinn an Wartbarkeit für den
    Mandanten überwiegt diesen Unterschied.
  \item \textbf{\ac{OIDC} und der duale \ac{URL}-Raum.} Dienste mit \ac{OIDC}
    funktionieren nicht gleichzeitig über die interne Domain und über Tor,
    weil \ac{OIDC} die Tokens fest an eine Basis-\ac{URL} bindet. Die
    Mechanik dahinter wurde im Abschnitt~\ref{sec:Tor} erläutert. Es handelt sich um
    eine reine Komforteinschränkung und nicht um ein Sicherheitsproblem. Als
    Komfortverbesserung wird dem Mandanten ergänzend ein WireGuard-Zugang
    empfohlen. Damit sind die internen Adressen von überall über denselben
    Pfad erreichbar, der \ac{URL}-Wechsel entfällt und die
    \ac{OIDC}-Anmeldung funktioniert ortsunabhängig. Tor bleibt als
    zusätzlicher anonymer Zugangsweg davon unberührt.
  \item \textbf{Collabora.} Der Office-Dienst benötigt für das Rendern von
    Dokumenten erweiterte Rechte, unter anderem die Capability
    \texttt{SYS\_ADMIN} sowie deaktivierte Seccomp- und AppArmor-Profile.
    Unter den für den Mandanten erreichbaren, nutzerdatenverarbeitenden
    Diensten ist das der am wenigsten eingeschränkte Container im Stack. Die
    in Abschnitt~\ref{sec:apphost-container} genannten Ausnahmen bei
    Monitoring und Docker-Anbindung betreffen dagegen reine
    Infrastrukturkomponenten ohne direkten Netzwerkzugriff von außen. Dieser
    Zustand wird akzeptiert, weil die Funktion sonst nicht läuft. Das Risiko
    wird über die Netzwerksegmentierung und die Anmeldung an OpenCloud
    begrenzt. Aus demselben Grund ist Collabora standardmäßig nicht über Tor
    erreichbar. In den letzten Monaten hat sich \ac{EU}-Office als
    (\ac{EU}-)Fork von OnlyOffice etabliert. Die Wahl fiel dennoch bewusst
    auf Collabora, weil es in der Open-Source-Community etabliert ist und
    eine große Nutzerbasis hat. \ac{EU}-Office ist noch sehr jung und
    erreicht bislang nicht dieselbe Reife und Stabilität. Die Entwicklung
    wird beobachtet und ein Umstieg bleibt für die Zukunft möglich.
  \item \textbf{Kata Containers und gVisor.} Beide zusätzlichen
    Sandbox"=Runtimes sind vorbereitet, aber standardmäßig nicht aktiv (siehe
    Abschnitt~\ref{sec:kata_konfig}).
\end{itemize}

\subsection{Backups und Wiederherstellung}

Für die Sicherung wird bewusst die in Proxmox integrierte Backup-Funktion
genutzt. Proxmox sichert komplette \acp{VM} als konsistente Snapshots und
kann diese regelmäßig und vollautomatisch wegschreiben. Ausschlaggebend für
diese Entscheidung ist die Einfachheit für den Mandanten. Da der gesamte
Apphost in einer einzigen \ac{VM} liegt, sichert ein \ac{VM}-Backup
automatisch den konsistenten Gesamtzustand aller Dienste auf einmal. Es kann
somit keine Situation entstehen, in der nach einer Wiederherstellung die
Datenbank eines Dienstes nicht mehr zu dessen Dateien passt. Gerade für einen
Anwender ohne tiefes \ac{IT}-Wissen ist es ein großer Vorteil, mit einem
Klick wieder einen funktionierenden Stand zu erhalten.

Damit entsteht eine gewisse Nähe zu Proxmox. Dieser Lock-in ist jedoch
gering. Das Backup-Format ist offen und gut dokumentiert. Die Sicherungen
sind durch Kompression und Deduplizierung effizient und lassen sich
verschlüsseln. Die Daten sind dadurch auch dann geschützt, wenn ein
Backup-Medium gestohlen wird. Für die Wiederherstellung nach einem
Hardwaredefekt genügt es, Proxmox neu aufzusetzen und die \ac{VM} aus dem
letzten Backup zurückzuspielen.

\subsubsection{Dateisystemintegrität mit \acs{AIDE}}
\ac{AIDE} prüft die Integrität des
Dateisystems. Dazu legt es Prüfsummen und Metadaten (Rechte, Eigentümer,
Timestamps) von zuvor konfigurierten Dateien und Verzeichnissen in einer
Datenbank ab und vergleicht sie bei jedem weiteren Lauf mit dem aktuellen
Zustand.

Ziel ist es, unautorisierte Veränderungen schnell sichtbar zu machen.
Beispiele sind eine eingeschleuste Backdoor in einer Binary, eine
manipulierte Konfigurationsdatei oder leicht zu übersehende Änderungen wie
veränderte \ac{SUID}-Bits. Zugrunde liegt die Annahme, dass ein Angreifer
einige Schutzmechanismen bereits überwunden und Spuren seiner
Kompromittierung auf dem Dateisystem hinterlassen hat. Diese Spuren lassen
sich über \ac{AIDE} auffinden.
\ac{AIDE} läuft automatisiert über einen systemd-Timer, das moderne
Äquivalent zu einem Cronjob. Die Ergebnisse müssen jedoch manuell geprüft
werden.

\subsection{Sicherheitshärtung im Überblick}

Das System kombiniert eine große Zahl einzelner Härtungsmaßnahmen, die alle
deklarativ in den NixOS-Modulen festgelegt sind. Hier folgt ein geordneter
Überblick, während die vollständige Liste mit konkreten Werten im Anhang
steht.

\begin{itemize}
  \item \textbf{Boot und Firmware.} Zum Einsatz kommen Secure Boot per
    lanzaboote, ein eingebundenes \ac{TPM}, ein deaktivierter
    Bootloader-Editor, erzwungene Kernel-Modul-Signaturen und der
    Kernel-Lockdown.
  \item \textbf{Kernel.} Aktiviert sind die \ac{CPU}-Mitigationen (Spectre,
    Meltdown, MDS und weitere), vollständiges \ac{ASLR}, init-on-alloc und
    init-on-free, eine erzwungene \ac{IOMMU} und ein eingeschränkter Zugriff
    auf Kernel-Informationen. Hinzu kommt eine an den \ac{CIS}-Benchmark
    angelehnte Modul-Blacklist.
  \item \textbf{Netzwerk.} nftables arbeitet mit Default-Deny sowie
    rate-limitiertem \ac{SSH} und \ac{ICMP}, ergänzt um gehärtete
    Sysctl-Einstellungen. Zusätzlich sind \ac{DNS}-over-\ac{TLS} mit
    \ac{DNSSEC} und authentifizierte Zeit über \ac{NTS} aktiv.
  \item \textbf{Zugriff.} \ac{SSH} funktioniert ausschließlich
    schlüsselbasiert mit modernen und post-quanten-fähigen Verfahren.
    Ergänzend existieren ein gehärtetes sudo mit I/O-Logging, eine
    unveränderliche Benutzerdatenbank, ein gesperrtes Root-Konto und AppArmor
    als Mandatory Access Control.
  \item \textbf{Audit und Intrusion Detection.} Der Linux Audit Daemon
    erfasst sicherheitsrelevante Ereignisse. Fail2ban arbeitet mit
    progressivem Banning und \ac{AIDE} prüft täglich die
    Dateisystemintegrität.
  \item \textbf{Container.} Abgesichert wird über User-Namespace-Remapping,
    den nur über einen eingeschränkten Socket-Proxy möglichen Zugriff auf die
    Docker-\ac{API}, entfernte Capabilities, read-only Dateisysteme wo
    möglich und einen wöchentlichen \ac{CVE}-Scan.
  \item \textbf{Updates.} Das System aktualisiert sich täglich über
    NixOS-Auto-Upgrade. RenovateBot pflegt die Container inklusive
    Digest-Pinning für reproduzierbare Image-Referenzen
    (siehe Abschnitt~\ref{sec:apphost-digest-pinning}) und der Nix-Store wird
    regelmäßig aufgeräumt.
\end{itemize}

\subsection{Fazit}
Der Apphost bildet den zentralen Einstiegspunkt für den Mandanten und das
Rückgrat für sämtliche Dienste. Auf ihm laufen nicht nur die Alltags- und
Produktivanwendungen, sondern auch der Hot-Wallet-Stack des
\ac{BTC}-Zahlungsverkehrs (siehe Abschnitt~\ref{Krypto}). Damit vereint eine einzige,
deklarativ beschriebene und umfassend gehärtete Maschine die gesamte für den
Mandanten sichtbare Infrastruktur. Die durchgängige Beschreibung als Code
macht den Aufbau vor der Ausführung prüfbar und nach einem Ausfall
reproduzierbar wiederherstellbar. Aus Sicht des Mandanten entsteht so ein
System, das sich an einer Stelle betreiben, überwachen und warten lässt, ohne
dass die Absicherung der einzelnen Dienste darunter leidet.